\begingroup

\section{Complete periodized arithmetic source and finite-cutoff proofs}


\subsection{Periodized source control: complete comparison calculations}
\label{sec:periodized-source-intake-proofs}

This companion retains the complete supplied periodized-source note,
including P1--P64 and P26a--P26b, as its proof source. It supplies the
source regularity, coordinate and image calculations used in reading
those formulas, makes the finite Fourier projection part of the exact
operator diagram, and retains the sharper one-sided error furnished by
the omitted positive Gram. The adjoining density companion supplies the
explicit definitions and full exponential-tail and completion proof.
The original source, remainder basis, arithmetic unit, and relations are
retained throughout.

The asymmetric cutoff calculation below is also the calculation in
FC5--FC7 and the four-endpoint paragraph following FC13 of the adjoining
periodized-curvature control bridge. PSA23--PSA25 record its source
and determinant steps, with the later AW/SP spectral and actual-overlap
control propagated through the same original and sampled coefficient maps, and PSA27--PSA29 record the same finite Fourier
boundary diagram as FC19--FC20. These
are audit derivations of the same mathematics, not additional results
created by repeating the calculation.

Fix the original packet polynomial $h$, its full canceled zero orders,
and $g=2\xi$. Write $v_h=g/h$, with removable values understood, and
\[
 w_h(t)=\frac{1}{2\pi}|v_h(1/2+it)|^2,
 \qquad m_{h,k}=w_h^{*k}.
 \tag{PSA1}
\]
These definitions must be inserted before the first use of $v_h,w_h$ in
the supplied standalone note. They specify the original density; no
rescaling of $g$, $h$, or the tensor-sum coordinate is performed. The
full proof that these functions have the required exponential tails,
including canceled-root values, is in the adjoining density companion.

\subsubsection{The original coordinates and the complete source domain}

Let $k\ge1$ and let $F$ be a finite linear combination of tensors of
strong-Schwartz functions on $(0,\infty)$ and their logarithmic
derivatives. Use $\mathcal K=\mathbb C$ for $k=1$ and
$\mathcal K=L^2(\mathbb R^{k-1},d\mathbf z)$ otherwise. In the ordered
coordinates $(r,z_1,\ldots,z_{k-1})$, put
\[
 y_i=r+z_i\ (i<k),\quad y_k=r,\quad x_i=e^{y_i}.
 \tag{PSA2}
\]
The derivative matrix from $(r,\mathbf z)$ to $(y_1,\ldots,y_k)$ has
first column all ones, last row $(1,0,\ldots,0)$ and the identity in the
remaining $k-1$ columns. Expansion along its last row gives determinant
$(-1)^{k-1}$. Thus the positive density is
$d^kx=e^{kr+\sum z_i}\,dr\,d\mathbf z$, with the orientation sign
recorded separately. The map and its inverse are
\[
 \begin{split}
 (\mathcal U_k F)(r,\mathbf z)
 &=e^{(kr+\sum z_i)/2}F(e^{r+z_1},\ldots,e^{r+z_{k-1}},e^r),\\
 (\mathcal U_k^{-1}\psi)(\mathbf x)
 &=(x_1\cdots x_k)^{-1/2}
 \psi(\log x_k,\log x_1-\log x_k,\ldots,
                  \log x_{k-1}-\log x_k).
 \end{split}
 \tag{PSA3}
\]
Substitution into the integral proves equality of the two squared
norms and proves that these maps are mutually inverse Hilbert
isometries. With $D_i=-x_i\partial_{x_i}$ and
$D^{(k)}=\sum_iD_i$, simultaneous change of all $y_i$ is exactly
change of $r$ with every $z_i$ fixed. Differentiating the displayed
factor, without altering the coordinate, gives
\[
 \mathcal U_kD^{(k)}=(-\partial_r+k/2)\mathcal U_k,
 \qquad
 \partial_r^j\mathcal U_k
 =(-1)^j\mathcal U_k(D^{(k)}-k/2)^j.
 \tag{PSA4}
\]

For every $a>0$ and integer $j\ge0$, strong-Schwartz estimates on each
original tensor factor imply
\[
 I_{a,j}:=\int_{\mathbb R} e^{2a|r|}
       \|\partial_r^j\mathcal U_kF(r)\|_{\mathcal K}^2dr<\infty.
 \tag{PSA5}
\]
Indeed after (PSA3)--(PSA4) this is the original $d^kx$ integral of
$|(D^{(k)}-k/2)^jF|^2$ multiplied by
$\max\{x_k^{2a},x_k^{-2a}\}$. The maximum is bounded by the sum of
those two weights, and every term is a finite integral of a finite sum
of products of original strong-Schwartz functions. This proves (PSA5)
on the stated source without asserting it for arbitrary Hilbert inputs.

For completeness it also yields normal convergence of periodization.
For any Hilbert-valued $H^1$ function $f$ on $[r,r+1]$, select a point
where its squared norm is no larger than its integral, use the
fundamental theorem of calculus, and apply Cauchy--Schwarz. The result is
$\|f(r)\|^2\le2\int_r^{r+1}(\|f(t)\|^2+\|f'(t)\|^2)dt$.
Consequently, for $\psi=\mathcal U_kF$,
\[
 \|\partial_r^j\psi(r)\|_{\mathcal K}^2
 \le 2e^{2a}e^{-2a|r|}(I_{a,j}+I_{a,j+1}).
 \tag{PSA6}
\]
The bound is summable over $r+mL$, uniformly for $r$ in a fixed
period and for every derivative. Thus the periodization is smooth;
the ensuing Fourier transforms, products, and unfolded correlations
are absolutely convergent on each finite source.

\subsubsection{The full circle image and its Fourier constants}

Use the plus-sign transform
$\widehat\psi(u)=\int\psi(r)e^{iur}dr$ and the orthonormal circle
basis $e_n(r)=L^{-1/2}e^{-2\pi inr/L}$. Unfolding its coefficient gives
\[
 \langle e_n,\mathcal P_L\psi\rangle
 =L^{-1/2}\int_{\mathbb R}\psi(r)e^{2\pi inr/L}dr,
 \qquad \mathcal P_L\psi=\sum_m\psi(r+mL).
 \tag{PSA7}
\]
Here the scalar basis is paired with the vector-valued function using
the conjugate-linear first-slot convention. In particular the zero
coefficient is $L^{-1/2}\widehat\psi(0)$, whereas the constant circle
function is $L^{-1}\widehat\psi(0)$.

The augmented coefficient map sends $\psi$ to the nonzero coefficients
and the unscaled vector $\widehat\psi(0)$. Its last summand has metric
$L^{-1}\langle\ ,\ \rangle_{\mathcal K}$. The inverse coordinate map
on this image inserts $L^{-1/2}\widehat\psi(0)$ at $n=0$. Thus, for a
finite column map $\Psi=\mathcal U_kT$ and $Z_0c=\widehat{\Psi c}(0)$,
\[
 M(L)=M^\times(L)+L^{-1}Z_0^*Z_0.
 \tag{PSA8}
\]
For $k=1$, $T=\Theta\phi$, the Connes--Consani map remains
$\mathcal E\phi=\tfrac12\mathcal U_1\Theta\phi$. It therefore has
zero coefficient $\mathcal M\Theta\phi(1/2)/(2\sqrt L)$. For the
original seed this is $\xi(1/2)/\sqrt L$. Its entire Gram is one quarter
of the original theta-source Gram. The coefficient and metric formulas
record that factor explicitly.

The scalar convolution formula retains the relative Hilbert fibre.
For $\psi=\mathcal U_kF_h^{\otimes k}$, take the further plus-sign
Fourier transform in $\mathbf z$. Its phase becomes
\[
 ur+\sum_{i<k}t_i z_i
 =\sum_{i<k}t_i y_i+
       \left(u-\sum_{i<k}t_i\right)y_k.
 \tag{PSA9}
\]
Fubini therefore gives the full amplitude
$\prod_{i<k}v_h(1/2+it_i)\,
 v_h(1/2+i(u-\sum_{i<k}t_i))$.
Plancherel in exactly the $k-1$ relative variables gives
\[
 \begin{split}
 \|\widehat\psi(u)\|_{\mathcal K}^2
 &=(2\pi)^{-(k-1)}\int_{\mathbb R^{k-1}}
   \prod_{i<k}|v_h(1/2+it_i)|^2
   |v_h(1/2+i(u-\sum_{i<k}t_i))|^2\,d\mathbf t\\
 &=2\pi m_{h,k}(u).
 \end{split}
 \tag{PSA10}
\]
By (PSA4), a polynomial column contributes its factor
$P(k/2+iu)$. Parseval and (PSA7) now prove the entire sampled matrix
\[
 M_N(L)_{ab}=\frac{2\pi}{L}\sum_{n\in\mathbb Z}
 m_{h,k}(2\pi n/L)
 \overline{P_a(k/2+2\pi in/L)}P_b(k/2+2\pi in/L).
 \tag{PSA11}
\]
The density companion proves positive weights for $k\ge2$ and hence
strict positivity of this finite polynomial Gram for every $L>0$.
For $k=1$, quantitative source comparison below supplies positivity
once the stated period has been chosen; no assertion about positivity
at every integer atom is substituted for that comparison.

\subsubsection{The original-source correlation inequality}

Define $M=\Psi^*\Psi$ and
$M_{a,\pm}=(e^{\pm ar}\Psi)^*(e^{\pm ar}\Psi)$ in the unchanged
coefficient basis. In original coordinates these are exactly the
integrals weighted by $x_k^{\pm2a}$. With
$c^*\mathcal C(s)d=\int\langle\Psi c(r),\Psi d(r+s)\rangle dr$,
unfolding pairs of translates gives
\[
 M(L)=\sum_{m\in\mathbb Z}\mathcal C(mL),\quad
 \mathcal C(0)=M,\quad\mathcal C(-s)=\mathcal C(s)^*.
 \tag{PSA12}
\]
For $s>0$, insert $e^{-ar}$ on the first factor and $e^{a(r+s)}$
on the second. Their product is $e^{as}$ times the original integrand.
Cauchy--Schwarz and the substitution $t=r+s$ give
\[
 |c^*\mathcal C(s)c|
 \le e^{-as}(c^*M_{a,-}c)^{1/2}(c^*M_{a,+}c)^{1/2}.
 \tag{PSA13}
\]
The roles of the weights reverse for negative $s$. The two geometric
series sum to $2/(e^{aL}-1)$. Applying
$2\sqrt{xy}\le x+y$ to these actual two quadratic forms yields
\[
 -\frac{M_{a,+}+M_{a,-}}{e^{aL}-1}
 \preceq M(L)-M\preceq
 \frac{M_{a,+}+M_{a,-}}{e^{aL}-1}.
 \tag{PSA14}
\]
This proves an inequality on every coefficient vector, including
complex cross-pairings, and hence the claimed matrix order.

\subsubsection{All-period recovery before the quotient minimum}

For $\delta>0$, quadratic-form Tonelli and
$u=2\pi n/L$ for $n>0$ give
\[
 \int_0^\infty L^{-2-\delta}
       \|\widehat{\Psi c}(2\pi n/L)\|^2dL
 =(2\pi n)^{-1-\delta}\int_0^\infty
       u^\delta\|\widehat{\Psi c}(u)\|^2du.
 \tag{PSA15}
\]
Negative $n$ supplies the negative half-line. The sum of
$n^{-1-\delta}$ over the positive integers is exactly
$\zeta(1+\delta)$, with no double counting of either half-line.
Equality of quadratic forms proves
\[
 M_\delta:=\frac1{2\pi}\int|u|^\delta
       \widehat\Psi(u)^*\widehat\Psi(u)du
 =\frac{(2\pi)^\delta}{\zeta(1+\delta)}
       \int_0^\infty L^{-1-\delta}M^\times(L)dL.
 \tag{PSA16}
\]
For $0<\delta\le\delta_0$, the mean-value formula implies
$||u|^\delta-1|\le\delta|\log|u||
(1+|u|^{\delta_0})$. The right side is integrable against every
finite-source Fourier Gram: the logarithm is locally integrable and
the tails are Schwartz. Thus the exact matrix with this logarithmic
integrand, denoted $\mathcal L_{\delta_0}$ in P31, obeys
$-\delta\mathcal L_{\delta_0}\preceq M_\delta-M
\preceq\delta\mathcal L_{\delta_0}$, proving convergence and its rate.
If the nonzero $Z_0^*Z_0/L$ term were inserted into (PSA16), its
small-period integral would be $\int_0 L^{-2-\delta}dL$ times that
matrix and would diverge. Formula (PSA8) is the exact reconstruction
of the omitted coefficient. Formula (PSA16) averages source Grams;
it makes no interchange of this integral with inverse compression.

\subsubsection{The same finite quotient and its actual relation coefficient}

Let $E=\mathbb C[S]/(\chi)$, $q=\deg\chi>0$, and
$\mathcal P_N=\{P:\deg P\le N\}$, with its original basis.
For $N\ge q-1$, the monic remainder map $J_N:\mathcal P_N\to E$
is onto. The actual arithmetic observation in the larger tensor target
is still $\eta J_N$, with
$\eta[P]=\upsilon_h^{\otimes k}P(A_k)1$ and
$\upsilon_h=j_h(g/h)$. We take determinants only in the stated
remainder basis of $E$.

For any positive source Gram $M_N$, set
\[
 K_N=J_NM_N^{-1}J_N^*,\quad G_N=K_N^{-1},\quad
 C_N=M_N^{-1}J_N^*G_N.
 \tag{PSA17}
\]
Surjectivity makes $K_N$ positive definite. Direct multiplication
gives $J_NC_N=I$. If $J_NP=v$, then $P=C_Nv+z$ with $J_Nz=0$.
The cross-term vanishes because
$C_N^*M_Nz=G_NJ_Nz=0$. Therefore
\[
 P^*M_NP=v^*G_Nv+z^*M_Nz,
 \qquad v^*G_Nv=\min_{J_NP=v}P^*M_NP.
 \tag{PSA18}
\]
This proves both the actual minimum and its unique representative.
Applying any source lower and upper bounds to this same affine set
and then taking its minimum transfers those same factors to $G_N$.

In particular, (PSA14) and the actual Rayleigh quotient
$\kappa_{a,N}=\max_{c\ne0}
c^*(M_{a,+,N}+M_{a,-,N})c/(c^*M_Nc)$ prove P34--P37.
The constraint set and coefficient coordinates are unchanged; the
minimizing representatives may change and their difference is retained
below. For the second representative
$C_N(L)$ obtained from $M_N(L)$,
\[
 J_N(C_N(L)-C_N)=0,\qquad
 C_N(L)-C_N=M_\chi\operatorname{div}_\chi(C_N(L)-C_N).
 \tag{PSA19}
\]
Monic division here is applied column by column and yields the actual
degree-$\le N-q$ coefficient, with the zero coefficient if $N=q-1$.
The tensor relation is also explicit. In the original polynomial
ring, fixed-order monic division in the variables $s_1,\ldots,s_k$
expresses $\chi(\sum s_i)Q(\sum s_i)$ as $\sum_i h(s_i)Q_i$.
Its remainder is zero because $\chi$ is the annihilator of the
original tensor-sum cyclic action. Replace the $i$th $F_h$ by
$\phi_*$, apply $Q_i(D_1,\ldots,D_k)$, and attach
$(-1)^{i-1}$. The tensor differential supplies the other
$(-1)^{i-1}$; their product is $+1$, and
$h(D_i)F_h=\Theta\phi_*$ supplies the original relation. Its
arithmetic image is the receiving supported zero, with its source
coefficient retained, and not external absence.

\subsubsection{The finite frequency cutoff and its one-sided allowance}

For $p\ge1$, define the unchanged-source derivative Gram
\[
 H_{p,N}=\int(1+r^2)(\partial_r^p\Psi_N)^*
                            \partial_r^p\Psi_N\,dr.
 \tag{PSA20}
\]
By (PSA4), the original-coordinate source map is
$(-1)^p(D^{(k)}-k/2)^p$ and its weight is
$1+(\log x_k)^2$. Integration by parts has no boundary term by
(PSA6), and gives
$\widehat{\partial_r^p\psi}(u)=(-iu)^p\widehat\psi(u)$.
Weighted Cauchy--Schwarz with
$\int(1+r^2)^{-1}dr=\pi$ proves
$\|\widehat{\Psi_Nc}(u)\|^2\le\pi|u|^{-2p}c^*H_{p,N}c$
for $u\ne0$.

Let $\Pi_{L,J}$ be the orthogonal projection onto precisely the
Fourier modes $|n|\le J$, including zero, and form
$M_N(L,J)=(\Pi_{L,J}\mathcal P_L\Psi_N)^*
                (\Pi_{L,J}\mathcal P_L\Psi_N)$.
The omitted Gram is positive. Summing both signs of $n$ and using
$\sum_{n>J}n^{-2p}\le\int_J^\infty t^{-2p}dt$ proves
\[
 0\preceq M_N(L)-M_N(L,J)
 \preceq\frac1{2p-1}
             \left(\frac{L}{2\pi J}\right)^{2p-1}H_{p,N}.
 \tag{PSA21}
\]
The exact coefficient before the series was
$(2\pi/L)(L/(2\pi))^{2p}$; (PSA21) retains the two-sided lattice
factor and the original Plancherel convention.

Use the actual polynomial inclusion from degree $N$ into degree $D$.
It pulls back each of $M_D,M_{a,\pm,D},H_{p,D}$ to the corresponding
degree-$N$ form. Consequently their Rayleigh bounds at $D$ control all
$N\le D$. With the exact P62 quantities, write
\[
 \eta_{\rm per}=\frac{\kappa_{a,D}}{e^{aL}-1},\qquad
 \eta_{\rm tail}=\frac{\lambda_{p,D}}{2p-1}
          \left(\frac{L}{2\pi J}\right)^{2p-1},\qquad
 \lambda_{p,D}=\max_{c\ne0}\frac{c^*H_{p,D}c}{c^*M_Dc}.
 \tag{PSA22}
\]
Choose $L$ and integer $J$ by P64 for specified positive budgets
whose sum is below one. These choices prove that
$\eta_{\rm per}+\eta_{\rm tail}<1$ after choosing the budgets
strictly below one. Combining the two actual inequalities gives the
sharper, asymmetric result
\[
 (1-\eta_{\rm per}-\eta_{\rm tail})M_N
 \preceq M_N(L,J)\preceq(1+\eta_{\rm per})M_N.
 \tag{PSA23}
\]
The upper bound retains the positivity of the omitted Gram. It implies
the supplied symmetric total-error bound, and also proves strict
positivity of the finite sampled matrix. For $k\ge2$, $J\ge D$ in
P64 separately supplies more distinct positive sample atoms than the
degree of a nonzero polynomial.

Define $G_N(L,J),C_N(L,J)$ by (PSA17) using this finite Gram.
By (PSA18), the factors in (PSA23) pass unchanged to $G_N(L,J)$.
The same proof of (PSA19) supplies its original theta correction.
Taking determinants on the fixed $q$-dimensional remainder space gives
\[
 q\log(1-\eta_{\rm per}-\eta_{\rm tail})
 \le\log\frac{\det G_N(L,J)}{\det G_N}
 \le q\log(1+\eta_{\rm per}).
 \tag{PSA24}
\]
For $D=2q+2$, all four endpoint degrees and both generator raises
remain among these actual restrictions. The common space for the
four determinant terms themselves is $\mathcal H=\mathcal P_{2q}$.
Let $I_N:\mathcal P_N\hookrightarrow\mathcal H$ and
$I_{2q,D}:\mathcal H\hookrightarrow\mathcal P_D$ be the literal
coefficient inclusions. Restriction gives
\[
 M^{(0)}=I_{2q,D}^*M_DI_{2q,D}=M_{2q},\qquad
 M^{(1)}=I_{2q,D}^*M_D(L,J)I_{2q,D}=M_{2q}(L,J).
 \tag{PSA24a}
\]
Set $\alpha=1-\eta_{\rm per}-\eta_{\rm tail}>0$ and
$\beta=1+\eta_{\rm per}$. These two forms satisfy
$\alpha M^{(0)}\preceq M^{(1)}\preceq\beta M^{(0)}$ by
(PSA23). In this paragraph let $x\in[0,1]$ be a metric path
parameter, distinct from the physical $x_i$, and retain
\[
 M(x)=(1-x)M^{(0)}+xM^{(1)},\quad
 D_{L,J}=(M^{(0)})^{-1}M^{(1)},\quad
 C(x)=M(x)^{-1}(M^{(1)}-M^{(0)}).
 \tag{PSA24b}
\]
These are the original and actual sampled forms, rather than a
Gamma comparison. The finite algebra of the AW and SP calculations
applies through the identity of this coefficient space, with their
endpoint forms replaced by exactly $M^{(0)},M^{(1)}$.
We prove that algebra for this source here.

For $N\in\{q-1,q,2q-1,2q\}$ put
$B_N=\times\chi:\mathcal P_{N-q}\to\mathcal P_N$ and define
\[
 \begin{aligned}
 P_N(x)&=I_N(I_N^*M(x)I_N)^{-1}I_N^*M(x),\\
 Q_N(x)&=I_NB_N(B_N^*I_N^*M(x)I_NB_N)^{-1}
                             B_N^*I_N^*M(x),\\
 U(x)&=Q_{2q-1}+Q_{2q}-Q_q,\qquad Q_{q-1}=0,\\
 W(x)&=(P_{2q-1}-P_{q-1})+(P_{2q}-P_q).
 \end{aligned}
 \tag{PSA24c}
\]
Every inverse is on the displayed actual domain. For $N=q-1$
the relation domain is zero and its determinant is one.
Multiplication verifies that these are the orthogonal projections
onto the specified polynomial and relation subspaces in $M(x)$.
The degree-less-than-$q$ remainder section $s_N$ and the monic
relation columns give a frame $[s_N,B_N]$ of determinant one:
its high-degree block is triangular with diagonal one.
Block elimination in its Gram therefore gives
\[
 V_N(x)=\frac{\det(I_N^*M(x)I_N)}
                  {\det(B_N^*I_N^*M(x)I_NB_N)},\quad
 \frac{d}{dx}\log V_N(x)=\operatorname{Tr}((P_N-Q_N)C(x)).
 \tag{PSA24d}
\]
This determinant is the same quotient determinant as (PSA17)--(PSA18),
since its Schur complement is the minimum Gram over the same
remainder fibres. The derivative follows by the finite determinant
derivative and cyclicity of trace. In the relation-first orientation,
the unchanged frame has determinant $(-1)^{q(N-q+1)}$ and its
squared modulus is one, so no orientation factor has been removed.
Consequently the four original signs give
$\mathcal B'(x)=\operatorname{Tr}((U-W)C(x))$.

The flags of relation subspaces have dimensions $1,q,q+1$, and
the polynomial flag has dimensions $q,q+1,2q,2q+1$.
Nested projections act as identity or zero on successive
orthogonal pieces. Hence each of $U,W$ has eigenvalues
$0$ of multiplicity $q$, $1$ of multiplicity $2$, and $2$
of multiplicity $q-1$. Thus each has rank $q+1$, trace $2q$
and trace of the square $\operatorname{Tr}U^2=\operatorname{Tr}W^2=4q-2$.
The eigenvalue-two multiplicity is zero at $q=1$.
Let $s(x)=\dim(\operatorname{ran}U\cap\operatorname{ran}W)\ge1$,
and write $c_1(x)\le\cdots\le c_n(x)$ for the eigenvalues of
$C(x)$, $n=2q+1$, and $d(x)=c_n(x)-c_1(x)$. Define
\[
 \begin{aligned}
 S_{\rm AW}^{L,J}(x)&=c_{q+2}(x)-c_q(x)
       +2\sum_{j=1}^{q-1}(c_{2q+2-j}(x)-c_j(x)),\\
 S_{\rm SP}^{L,J}(x)&=\frac{d(x)}2\min\left\{
      4q-2s(x),\sqrt{(2q+1)(8q-4-2\operatorname{Tr}(UW))}\right\},\\
 \mathcal J_{L,J}&=\int_0^1
            \min\{S_{\rm AW}^{L,J}(x),S_{\rm SP}^{L,J}(x)\}\,dx.
 \end{aligned}
 \tag{PSA24e}
\]
For an orthogonal projection $P$ of rank $r$, its diagonal entries
in a $C(x)$-eigenbasis lie in $[0,1]$ and sum to $r$.
Moving a fixed diagonal mass from a smaller eigenvalue to a larger
one increases its trace pairing; hence
$\sum_{j=1}^rc_j\le\operatorname{Tr}(PC)
\le\sum_{j=n-r+1}^nc_j$.
Each of $U,W$ is the sum of a rank-$q+1$ range projection and a
rank-$q-1$ eigenvalue-two projection. Applying these bounds and
subtracting their full extrema yields
$|\operatorname{Tr}((U-W)C)|\le S_{\rm AW}^{L,J}$.

Let $E_L$ be the projection onto the actual intersection of the
two ranges. The positive operators $U-E_L,W-E_L$ both have trace
$2q-s(x)$, because each of $U,W$ dominates its range projection.
Thus $\|U-W\|_1\le4q-2s(x)$. Also
$\operatorname{Tr}(U-W)=0$ and
$\operatorname{Tr}((U-W)^2)=8q-4-2\operatorname{Tr}(UW)$.
Cauchy--Schwarz on its $n$ real eigenvalues gives the other
trace-norm estimate in (PSA24e).
Subtracting $(c_1+c_n)I/2$ from $C$ leaves the trace pairing
unchanged and bounds its absolute value by $d(x)\|U-W\|_1/2$,
as is seen in an orthonormal eigenbasis of $U-W$.
This proves the SP bound on the same signed derivative, and
therefore its pointwise minimum with the AW bound.
Smooth positive forms give bounded projector matrices and
continuous eigenvalues; rank strata are defined by matrix minors,
so $s(x)$ is Borel. The minimum is integrable, and integration
with its original signs gives the first bound below.


The actual source in (PSA24a) has additional moment structure.
The first measure is $m_{h,k}(y)\,dy$, and the second is the
unchanged positive atomic measure
$\sum_{|n|\le J}(2\pi/L)m_{h,k}(2\pi n/L)\delta_{2\pi n/L}$,
including zero. Their degree-$2q$ Grams are exactly
$M^{(0)},M^{(1)}$ by (PSA10)--(PSA11) and (PSA21).
Positivity was proved in (PSA23). Hence the complete RMT
calculation applies with this original/sampled pair, with no
change to any atom, mass or polynomial coefficient.
For $a_q=2q-1$ put
\[
 \begin{aligned}
 A_{L,J}(x)&=a_q\sqrt{1-e^{-\mathcal B(x)/a_q}}\,d(x),\\
 \mathcal J^{(3)}_{L,J}&=\int_0^1
 \min\{S_{\rm AW}^{L,J}(x),S_{\rm SP}^{L,J}(x),A_{L,J}(x)\}\,dx,\\
 \mathcal I_{L,J}&=\int_0^1\min\left\{d(x),
 \frac{S_{\rm AW}^{L,J}(x)}{a_q\sqrt{1-e^{-\mathcal B(x)/a_q}}},
 \frac{S_{\rm SP}^{L,J}(x)}{a_q\sqrt{1-e^{-\mathcal B(x)/a_q}}}\right\}\,dx.
 \end{aligned}
 \tag{PSA24f}
\]
RMT11 keeps both original $q$-angle lists of the crossed quotient
pairs and proves $|\mathcal B'|\le A_{L,J}$ on these measures.
Its strict positivity argument applies also to this atomic
endpoint: equality $\mathcal B=0$ would place $1$ orthogonal to
$\chi\mathcal P_q$, and pairing with
$\chi^{\#_k}(S)=\sum_j\overline{\chi_j}(k-S)^j$ would give
$0=\int|\chi|^2\,d\mu_x$, contradicting the positive Gram of
the nonzero $\chi$. Thus $\mathcal B>0$ on the whole compact path.
Taking the minimum of this angle derivative bound and the two
fully proved bounds above yields
$|\mathcal B(1)-\mathcal B(0)|\le\mathcal J^{(3)}_{L,J}
\le\mathcal J_{L,J}$. The generic two-form proof remains valid;
this further inequality uses precisely the actual moment measures.

Let $b_1\le\cdots\le b_n$ be all eigenvalues of $D_{L,J}$.
The form comparison gives $\alpha\le b_j\le\beta$.
Factoring $M(x)=M^{(0)}((1-x)I+xD_{L,J})$ proves
$c_j(x)=(b_j-1)/(1-x+xb_j)$; the order follows from derivative
$(1-x+xb_j)^{-2}>0$ with respect to $b_j$.
Its integral is $\log b_j$. Therefore the full strengthened
finite-circle error, with every original source constant, is
\[
 \begin{aligned}
 |\mathcal B_{h,k}(L,J)-\mathcal B_{h,k}|
 &\le\mathcal J^{(3)}_{L,J}\le\mathcal J_{L,J}\\
 &\le\log\frac{b_{q+2}}{b_q}
       +2\sum_{j=1}^{q-1}\log\frac{b_{2q+2-j}}{b_j}\\
 &\le(2q-1)\log\frac{1+\eta_{\rm per}}
                         {1-\eta_{\rm per}-\eta_{\rm tail}}.
 \end{aligned}
 \tag{PSA25}
\]
The first quantity is also at most $\int_0^1S_{\rm SP}^{L,J}(x)\,dx$.
Each logarithmic ratio in the middle expression is at most
$\log(\beta/\alpha)$ and its coefficient total is $2q-1$.
This proves the last inequality, including the empty sum at $q=1$.
It is at most $(2q-1)\log((1+\eta)/(1-\eta))$ for
$\eta=\eta_{\rm per}+\eta_{\rm tail}$. The older bound with
coefficient $2q$ follows from this stronger estimate; the
asymmetric omitted-tail allowance is unchanged.
Fixed positive budgets still give an $O(q_k)$ error, whose
ratio to $q_k\log k$ tends to zero. No bound on the growth of
the original $\kappa_{a,D}$ or $\lambda_{p,D}$ is inserted.
The sampled values remain the exact analytic values with every
zero mode, Fourier factor and tensor source coefficient retained.



This also improves the actual sampled endpoint itself in its
nonlinear coordinate. Define
$H_0=2\operatorname{arcosh}(e^{\mathcal B_{h,k}/(2a_q)})$.
RMT16's derivative and inverse, on this identified moment path,
give $|\mathscr H_{a_q}(\mathcal B_{h,k}(L,J))-H_0|
\le\mathcal I_{L,J}\le\log(b_n/b_1)\le\log(\beta/\alpha)$.
Combining the increasing inverse with the just-proved additive
bound yields
\[
 \begin{aligned}
 \mathcal B_{h,k}(L,J)&\ge\max\left\{0,
 \mathcal B_{h,k}-\mathcal J^{(3)}_{L,J},
 2a_q\log\cosh\frac{\max\{0,H_0-\mathcal I_{L,J}\}}2\right\},\\
 \mathcal B_{h,k}(L,J)&\le\min\left\{
 \mathcal B_{h,k}+\mathcal J^{(3)}_{L,J},
 2a_q\log\cosh\frac{H_0+\mathcal I_{L,J}}2\right\}.
 \end{aligned}
 \tag{PSA25a}
\]
All denominators in (PSA24f) are justified by the strict positivity
above, with no uniform lower bound in tensor order presumed.

\subsubsection{The finite observation map and both Laplacian boundary terms}

Put $\Phi_{N,L,J}=\Pi_{L,J}\mathcal P_L\mathcal U_k
\mathcal V_{h,k}|_{\mathcal P_N}$. Strict positivity in (PSA23)
makes it injective. The exact observation has the stated image as
domain:
\[
 \mathcal J_{N,L,J}:\operatorname{im}\Phi_{N,L,J}\longrightarrow E,
 \qquad
 \mathcal J_{N,L,J}=J_N\Phi_{N,L,J}^{-1}.
 \tag{PSA26}
\]
On that image, the inverse can equivalently be written
$(\Phi_{N,L,J}^*\Phi_{N,L,J})^{-1}\Phi_{N,L,J}^*$ in the unchanged
coefficient coordinates. Its kernel is
$\Phi_{N,L,J}(\ker J_N)$, so its quotient is exactly $E$ with the
metric (PSA17), retaining all jets. This construction uses finite
dimensionality and the proved lower source bound; it asserts no
bounded inverse on a larger completed source.

On the periodic $H^1$ domain let $D_L=-\partial_r+k/2$.
The adjoint is $D_L^*=\partial_r+k/2$, so $D_L^*+D_L=k$.
On periodic $H^2$, $\mathscr L_{L,k}=D_L^2-kD_L
=\partial_r^2-k^2/4$ is the self-adjoint Fourier diagonal with
values $-(2\pi n/L)^2-k^2/4$. The projection $\Pi_{L,J}$ commutes
with both derivatives on these domains.

Let $r=\mathcal V_{h,k}C_N(L,J)$ and use the original arithmetic
action $A=M_S$ on $E$. The finite representative and its actual
boundary are
\[
 R=\Pi_{L,J}\mathcal P_L\mathcal U_kr,\qquad
 B=\Pi_{L,J}\mathcal P_L\mathcal U_k(D^{(k)}r-rA).
 \tag{PSA27}
\]
The unperiodized difference has zero remainder at degree $N+1$ and
has the monic theta primitive already proved above. Commutation and
(PSA4) prove $D_LR-RA=B$ and $R^*R=G_N(L,J)$. Substituting
$RA=D_LR-B$ and using $D_L^*+D_L=k$ gives
\[
 A^*G_N(L,J)+G_N(L,J)A-kG_N(L,J)
 =-(R^*B+B^*R).
 \tag{PSA28}
\]
Applying $D_L$ once more to the first identity and subtracting $k$
times it gives, with every term retained,
\[
 \mathscr L_{L,k}R-R(A^2-kA)=(D_L-k)B+BA.
 \tag{PSA29}
\]
This proves the finite-cutoff extension of P46--P48. Setting
$\Pi_{L,J}=I$ gives the original full-period identities.

For $Av=\rho v$, (PSA29) is
$(\mathscr L_{L,k}-\rho(\rho-k))Rv=((D_L-k)B+BA)v$.
On the original Fourier diagonal, the squared norm on the left is
the sum of
$|-(2\pi n/L)^2-k^2/4-\rho(\rho-k)|^2\|(Rv)_n\|^2$.
Every coefficient is at least
$|\operatorname{Im}(\rho(\rho-k))|^2$, proving P49 also after the
finite projection. The conclusion is a lower bound on this actual
boundary, not an upper bound on the arithmetic endpoint volume.
The exact observation (PSA26), its one- and two-degree extensions,
and (PSA27)--(PSA29) supply the map from the circle calculation to
the arithmetic action.
For those inverse observation extensions, take the common controlled
source degree $D\ge N+2$; the endpoint choice $D=2q+2$ already has
this property for every endpoint $N\le2q$. The boundary identity
itself uses the displayed smooth columns and needs injectivity only
at its representative degree $N$.

\subsubsection{The subsequent completion and the scope of verification}

The fixed-circle limit keeps $L$ fixed and completes the increasing
polynomial source. The adjoining density proof supplies its typed
isometry to $\ell^2(\mathbb Z,\nu)$, polynomial density for the
original measure and its $|\chi|^2$ weight, and the exact quotient
map $[P]\mapsto(P(S_n))_{\chi(S_n)=0}$. Its kernel is
$(p_L)/(\chi)$, with $p_L$ the squarefree product of the sampled
roots, and the canonical Grams decrease to the sampled-value Gram.
That proof retains all multiplicities in $\chi$ before taking the
calculated kernel. The finite observation (PSA26) instead retains
the entire remainder space at the specified $N,L,J$. These are
related by the explicit quotient and completion arrows; an
unsampled class or a nonconstant sampled jet maps to a supported
zero under the latter arrow.

The complete supplied 16-method finite suite passed both normally
and under Python optimization, with identical result records.
Seven additional in-memory mutations each failed at one intended
mathematical assertion in both modes, with no execution errors.
These finite checks cover stated algebraic interfaces. The original
zero-mode command-line negative is an unconditional failure-path
control; the additional mutation removes the actual zero-mode Gram
from its original complex-matrix calculation. The original raw
checker and proof sources remain unchanged. The analytic proofs
above and the adjoining density proof remain written mathematics;
no Lean or interval certificate is inferred from those checks.


\endgroup
